\begin{definition}
  For a list of natural digits written with the units digit first, define its base-$b$ value
  recursively by $V_b([])=0$ and $V_b(d::D)=d+bV_b(D)$.

  \begin{lean}
    def digitValue (b : ℕ) : List ℕ → ℕ

      | [] => 0

      | d :: ds => d + b * digitValue b ds
  \end{lean}
\end{definition}

\begin{definition}
  A digit list is \emph{canonical} when every digit is less than $b$ and its last, most
  significant digit is nonzero.
  The empty list is the canonical representation of zero.

  \begin{lean}
    def CanonicalDigits (b : ℕ) : List ℕ → Prop

      | [] => True

      | d :: ds => d < b ∧ (ds = [] → d ≠ 0) ∧ CanonicalDigits b ds
  \end{lean}
\end{definition}

\begin{theorem}
  \label{ent:thm:existsunique-digits}
  For each natural number $n$ and base $b>1$, there is exactly one canonical digit list with value
  $n$.
  For positive $n$, this list is nonempty.

  \begin{lean}
    theorem existsUnique_digits (b : ℕ) (hb : 1 < b) (n : ℕ) :
        ∃! ds : List ℕ, CanonicalDigits b ds ∧ digitValue b ds = n
  \end{lean}
\end{theorem}

\begin{proof}
  Use strong induction on $n$.
  Zero has the empty expansion.
  For positive $n$, division by $b$ gives $n=bq+r$, where $0\le r<b$ and $0\le q<n$.
  Expand $q$ by induction and prepend $r$; if the tail is empty, positivity of $n$ ensures $r\ne0$.

  \begin{lean}
    have hex : ∀ n : ℕ, ∃ ds : List ℕ, CanonicalDigits b ds ∧ digitValue b ds = n := by
      intro n
      induction n using Nat.strong_induction_on with
      | h n ih =>
        by_cases hn : n = 0

        · exact ⟨[], trivial, by exact hn.symm⟩

        · obtain ⟨⟨q, r⟩, ⟨heq, hr⟩, _⟩ :=
            Chapter01.existsUnique_quotient_remainder n ⟨b, by omega⟩
          change (n : ℤ) = q * b + r at heq
          change r < b at hr

          have hb' : (1 : ℤ) < b := by exact_mod_cast hb

          have hr' : (r : ℤ) < b := by exact_mod_cast hr

          have hq : 0 ≤ q := by nlinarith

          have hnpos : (0 : ℤ) < n := by omega

          have hqn : q < n := by nlinarith

          have hlt : q.toNat < n := by omega

          obtain ⟨ds, hds, hval⟩ := ih q.toNat hlt

          refine ⟨r :: ds, ⟨hr, ?_, hds⟩, ?_⟩

          · intro hnil hrzero
            rw [hnil] at hval
            change 0 = q.toNat at hval

            have : q = 0 := by omega
            rw [this, hrzero] at heq
            simp only [zero_mul, Nat.cast_zero, add_zero] at heq
            exact hn (by exact_mod_cast heq)

          · change r + b * digitValue b ds = n
            rw [hval]

            have hc : (q.toNat : ℤ) = q := Int.toNat_of_nonneg hq
            exact_mod_cast (show (r : ℤ) + (b : ℤ) * q.toNat = n by
                rw [hc]
                nlinarith)
  \end{lean}

  The construction gives existence.
  For uniqueness, equality of two positional values gives equal quotient--remainder pairs on
  division by $b$: the first digits agree and the tails have equal values.
  Repeating this argument proves equality of the complete lists.

  \begin{lean}
    obtain ⟨ds, hd, hv⟩ := hex n

    exact ⟨ds, ⟨hd, hv⟩, fun es he => digitValue_injective b hb es ds he.1 hd (he.2.trans hv.symm)⟩
  \end{lean}
\end{proof}

\begin{theorem}
  \label{ent:thm:existsunique-binary-digits}
  Every natural number has a unique canonical base-two expansion.
  Since each digit is less than two, the digits are zero and one.

  \begin{lean}
    theorem existsUnique_binary_digits (n : ℕ) :
        ∃! ds : List ℕ, CanonicalDigits 2 ds ∧ digitValue 2 ds = n
  \end{lean}
\end{theorem}

\begin{proof}
  Specialize the positional representation theorem to base two.

  \begin{lean}
    existsUnique_digits 2 (by decide) n
  \end{lean}
\end{proof}

\begin{definition}
  For integer coefficients written in increasing order of degree, define $P_b([])=0$ and
  $P_b(d::D)=d+bP_b(D)$.
  Thus the same notation evaluates polynomials and positional digit lists.

  \begin{lean}
    def positionalValue (b : ℤ) : List ℤ → ℤ

      | [] => 0

      | d :: ds => d + b * positionalValue b ds
  \end{lean}
\end{definition}

\begin{theorem}
  \label{ent:thm:positionalvalue-modeq}
  An integer polynomial takes congruent values at congruent arguments: if $a\equiv b\pmod n$, then
  $P(a)\equiv P(b)\pmod n$.

  \begin{lean}
    theorem positionalValue_modEq (n a b : ℤ) (h : a ≡ b [ZMOD n]) (ds : List ℤ) :
        positionalValue a ds ≡ positionalValue b ds [ZMOD n]
  \end{lean}
\end{theorem}

\begin{proof}
  Induct on the coefficient list.
  The zero polynomial gives the base case.
  The recursive step adds a common constant coefficient to congruent products.

  \begin{lean}
    induction ds with
    | nil => exact modEq_refl n 0

    | cons d ds ih => exact modEq_add (modEq_refl n d) (modEq_mul h ih)
  \end{lean}
\end{proof}

\begin{definition}
  The alternating sum starts with the units digit: $A([])=0$ and $A(d::D)=d-A(D)$.

  \begin{lean}
    def alternatingSum : List ℤ → ℤ

      | [] => 0

      | d :: ds => d - alternatingSum ds
  \end{lean}
\end{definition}

\begin{theorem}
  \label{ent:thm:nine-dvd-iff-digit-sum}
  A decimal positional value is divisible by $9$ precisely when the sum of its digits is divisible
  by $9$.
  The identity also holds for arbitrary integer coefficients.

  \begin{lean}
    theorem nine_dvd_iff_digit_sum (ds : List ℤ) :
        9 ∣ positionalValue 10 ds ↔ 9 ∣ ds.sum
  \end{lean}
\end{theorem}

\begin{proof}
  Since $10\equiv1\pmod9$, substitute one for ten in the positional polynomial.
  Its value at one is the digit sum.

  \begin{lean}
    have h := positionalValue_modEq 9 10 1 (by decide) ds
    rw [positionalValue_one] at h
    exact dvd_iff_of_modEq h
  \end{lean}
\end{proof}

\begin{theorem}
  \label{ent:thm:three-dvd-iff-digit-sum}
  A decimal positional value is divisible by $3$ precisely when the sum of its digits is divisible
  by $3$.
  The identity also holds for arbitrary integer coefficients.

  \begin{lean}
    theorem three_dvd_iff_digit_sum (ds : List ℤ) :
        3 ∣ positionalValue 10 ds ↔ 3 ∣ ds.sum
  \end{lean}
\end{theorem}

\begin{proof}
  The same argument uses $10\equiv1\pmod3$.

  \begin{lean}
    have h := positionalValue_modEq 3 10 1 (by decide) ds
    rw [positionalValue_one] at h
    exact dvd_iff_of_modEq h
  \end{lean}
\end{proof}

\begin{theorem}
  \label{ent:thm:eleven-dvd-iff-alternating-sum}
  A decimal positional value is divisible by $11$ precisely when the alternating sum of its
  digits, starting with the units digit is divisible by $11$.
  The identity also holds for arbitrary integer coefficients.

  \begin{lean}
    theorem eleven_dvd_iff_alternating_sum (ds : List ℤ) :
        11 ∣ positionalValue 10 ds ↔ 11 ∣ alternatingSum ds
  \end{lean}
\end{theorem}

\begin{proof}
  Use $10\equiv-1\pmod{11}$.
  Evaluation at minus one is the alternating sum.

  \begin{lean}
    have h := positionalValue_modEq 11 10 (-1) (by decide) ds
    rw [positionalValue_neg_one] at h
    exact dvd_iff_of_modEq h
  \end{lean}
\end{proof}

\begin{theorem}
  \label{ent:thm:two-dvd-iff-last-digit}
  A nonempty decimal expansion is divisible by two exactly when its units digit is even.

  \begin{lean}
    theorem two_dvd_iff_last_digit (d : ℤ) (ds : List ℤ) :
        2 ∣ positionalValue 10 (d :: ds) ↔ 2 ∣ d
  \end{lean}
\end{theorem}

\begin{proof}
  All other terms contain a factor of ten, which is divisible by two.

  \begin{lean}
    dvd_positionalValue_cons_iff 2 10 d ds (by decide)
  \end{lean}
\end{proof}

\begin{theorem}
  \label{ent:thm:five-dvd-iff-last-digit}
  For a units digit $0\le d<10$, a decimal expansion is divisible by five exactly when $d=0$ or
  $d=5$.

  \begin{lean}
    theorem five_dvd_iff_last_digit (d : ℤ) (ds : List ℤ) (hd : 0 ≤ d ∧ d < 10) :
        5 ∣ positionalValue 10 (d :: ds) ↔ d = 0 ∨ d = 5
  \end{lean}
\end{theorem}

\begin{proof}
  Remove the terms divisible by ten.
  A multiple of five in the interval from zero to nine is zero or five.

  \begin{lean}
    rw [dvd_positionalValue_cons_iff 5 10 d ds (by decide)]
    constructor

    · rintro ⟨k, hk⟩
      omega

    · rintro (rfl | rfl) <;> decide
  \end{lean}
\end{proof}

\begin{theorem}
  \label{ent:thm:four-dvd-iff-last-two-digits}
  A decimal positional value with units digit $d$ and tens digit $e$ is divisible by four exactly
  when $d+10e$ is divisible by four.

  \begin{lean}
    theorem four_dvd_iff_last_two_digits (d e : ℤ) (ds : List ℤ) :
        4 ∣ positionalValue 10 (d :: e :: ds) ↔ 4 ∣ d + 10 * e
  \end{lean}
\end{theorem}

\begin{proof}
  The remaining contribution is one hundred times the tail value, hence is divisible by four.

  \begin{lean}
    apply dvd_iff_of_modEq
    exact (modEq_iff_dvd_sub _ _ _).2
      ⟨25 * positionalValue 10 ds, by
          simp only [positionalValue]
          ring⟩
  \end{lean}
\end{proof}

\begin{theorem}
  \label{ent:thm:detects-single-digit-error}
  If $\gcd(10,w)=1$ and $d,e$ are distinct decimal digits, then $wd\not\equiv we\pmod{10}$.

  \begin{lean}
    theorem detects_single_digit_error (w d e : ℤ) (hw : Chapter01.gcd 10 w = 1)
        (hd : 0 ≤ d ∧ d < 10) (he : 0 ≤ e ∧ e < 10) (hne : d ≠ e) :
        ¬w * d ≡ w * e [ZMOD 10]
  \end{lean}
\end{theorem}

\begin{proof}
  Otherwise coprime cancellation would give $10\mid d-e$.
  Since $-10<d-e<10$, this forces equality of the digits.

  \begin{lean}
    intro h

    obtain ⟨k, hk⟩ := (modEq_iff_dvd_sub _ _ _).1 (modEq_cancel hw h)
    omega
  \end{lean}
\end{proof}
